% This LaTeX was auto-generated from MATLAB code.
% To make changes, update the MATLAB code and export to LaTeX again.

\documentclass{article}

\usepackage[utf8]{inputenc}
\usepackage[T1]{fontenc}
\usepackage{lmodern}
\usepackage{graphicx}
\usepackage{color}
\usepackage{hyperref}
\usepackage{amsmath}
\usepackage{amsfonts}
\usepackage{epstopdf}
\usepackage[table]{xcolor}
\usepackage{matlab}
\usepackage[paperheight=795pt,paperwidth=614pt,top=72pt,bottom=72pt,right=72pt,left=72pt,heightrounded]{geometry}

\sloppy
\epstopdfsetup{outdir=./}
\graphicspath{ {./lab5_media/} }

\matlabmultipletitles

\begin{document}

\matlabtitle{Исследование типовых динамических звеньев}

\textbackslash{}section\{Исследование типовых динамических звеньев\}

В рамках данной лабораторной работы необходимо исследовать реальные объекты, сопоставим их с известными нам типовыми динамическими звеньями и получим их передаточные функции в стандартизированной форме (через коэффициенты усиления $K$ и постоянные времени $T$). Затем запишем аналитические выражения для временных (переходной и весовой) и частотных (АЧХ, ФЧХ, ЛАЧХ и ЛФЧХ) характеристик исследуемых звеньев. Проведем моделирование с целью получения графического представления всех перечисленных характеристик и сравним с теоретическими для данного типового звена.

\textbf{Вариант 27}

\textbackslash{}subsection\{ДПТ\}

Даны уравнения двигателя постоянного тока независимого возбуждения:

$J\dot{\omega} =M$, $M=k_{\text{m}} I$, $I=\frac{U+\varepsilon_{\text{i}} }{R}$, $\varepsilon_i =-k_{\text{e}} \omega$.

Возьмем из \textbf{Таблицы 1} инструкции по выполнению данной л/р значения, которые соответствуют нашему варианту, для следующих величин:

• $k_{\text{m}} =0.3824$ Н·м/А — конструктивная постоянная по моменту;

• $k_{\text{e}} =0.3824$ В·с — конструктивная постоянная по ЭДС;

• $J=0.0019$ кг·м\textasciicircum{}2 — момент инерции ротора;

• $R=4.6296$ Ом — активное сопротивление обмоток ротора.

Входом объекта считаем $U(t)$, а выходом $\omega (t)$.

\textbf{Найдём передаточную функцию и определим тип звена:}

\begin{par}
$$\begin{array}{l}
J\dot{\omega} =k_m \cdot \frac{U-k_e \omega }{R}=\frac{k_m }{R}U-\frac{k_e k_m }{R}\omega ,~~J\dot{\omega} +\frac{k_e k_m }{R}\omega =\frac{k_m }{R}U\\
\frac{JR}{k_e k_m }\dot{\omega} +\omega =\frac{1}{k_e }U
\end{array}$$
\end{par}

Передаточная функция: $W(s)=\frac{K}{Ts+1}$, где $K=\frac{1}{k_e },~~T=\frac{JR}{k_e k_m }$.

Типовое звено - \textit{реальное усилительное}.

\textbf{Запишем аналитические выражения для временных и частотных характеристик:}

\begin{par}
$$W(j\omega )=\frac{K}{T(j\omega )+1}=\frac{K-KT(j\omega )}{T^2 \omega^2 +1}=\frac{K}{T^2 \omega^2 +1}-\frac{KT\omega }{T^2 \omega^2 +1}j$$
\end{par}

Весовая: $w(t)={\mathcal{L}}^{-1} \lbrace W(s)\rbrace ={\mathcal{L}}^{-1} \left\lbrace \frac{K}{Ts+1}\right\rbrace ={\mathcal{L}}^{-1} \left\lbrace \frac{K}{T}\cdot \frac{1}{s+\frac{1}{T}}\right\rbrace =\frac{K}{T}e^{-\frac{t}{T}}$

Переходная: $h(t)={\mathcal{L}}^{-1} \left\lbrace \frac{W(s)}{s}\right\rbrace ={\mathcal{L}}^{-1} \left\lbrace \frac{K}{s(Ts+1)}\right\rbrace ={\mathcal{L}}^{-1} \left\lbrace \frac{K}{s}-\frac{K}{s+\frac{1}{T}}\right\rbrace =K\left(1-e^{-\frac{t}{T}} \right)$

АЧХ: 

\begin{par}
$$\begin{array}{l}
A(\omega )=|W(j\omega )|=\sqrt{{\text{Re}}^2 (W)+{\text{Im}}^2 (W)}=\sqrt{{\left(\frac{K}{T^2 \omega^2 +1}\right)}^2 +{\left(\frac{KT\omega }{T^2 \omega^2 +1}\right)}^2 }=\\
=\sqrt{\frac{K^2 +T^2 K^2 \omega^2 }{(T^2 \omega^2 +1)^2 }}=\frac{K}{\sqrt{T^2 \omega^2 +1}}=
\end{array}$$
\end{par}

ФЧХ:

\begin{par}
$$\varphi (\omega )=\arg W(j\omega )=\text{atan2}(\text{Im}(W),\text{Re}(W))=\text{atan2}\left(-\frac{KT\omega }{1+T^2 \omega^2 },\frac{K}{1+T^2 \omega^2 }\right)=$$
\end{par}



ЛАЧХ и ЛФЧХ находятся из формулы: $\lg (W(j\omega ))=\lg (A(\omega ))+j\lg (e)\phi (\omega )$.

ЛАЧХ: $L(\omega )=20\lg \left(A(\omega )\right)=20\lg K-20\lg \sqrt{T^2 \omega^2 +1}=$

ЛФЧХ: $\varphi (\omega )=\text{atan}\left(\frac{\text{Im}(W)}{\text{Re}(W)}\right)=-\text{atan}(T\omega )=$






\textbackslash{}subsection\{ДПТ 2.0\}

Даны уравнения двигателя постоянного тока независимого возбуждения:

$J\dot{\omega} =M$, $M=k_{\text{m}} I$, $I=\frac{U+\varepsilon }{R}$, $\varepsilon =\varepsilon_{\text{i}} +\varepsilon_{\text{s}}$, $\varepsilon_{\text{i}} =-k_{\text{e}} \omega$, $\varepsilon_{\text{s}} =-L\dot{I}$.

Возьмем из \textbf{Таблицы 1} инструкции по выполнению данной л/р значения, которые соответствуют нашему варианту, для следующих величин:

• $k_{\text{m}} =0.3824$ Н·м/А — конструктивная постоянная по моменту;

• $k_{\text{e}} =0.3824$ В·с — конструктивная постоянная по ЭДС;

• $J=0.0019$ кг·м\textasciicircum{}2 — момент инерции ротора;

• $R=4.6296$ Ом — активное сопротивление обмоток ротора.

• $L=1.1194$ Гн — индуктивность обмоток ротора.

Входом объекта считаем $U(t)$, а выходом $\omega (t)$.

\textbf{Найдём передаточную функцию и определим тип звена:}

\begin{par}
 $$\begin{array}{l}
J\dot{\omega} =k_m \cdot \frac{U+\varepsilon }{R}=k_m \cdot \frac{U-k_e \omega -L\dot{I} }{R},~~JR\dot{\omega} =k_m U-k_e k_m \omega -Lk_m \dot{I} \\
JL\ddot{\omega} +JR\dot{\omega} +k_e k_m \omega =k_m U,~~\frac{JL}{k_e k_m }\ddot{\omega} +\frac{JR}{k_e k_m }\dot{\omega} +\omega =\frac{1}{k_e }U
\end{array}$$
\end{par}

Передаточная функция: $W(s)=\frac{K}{T^2 s^2 +2\xi Ts+1}$, где $K=\frac{1}{k_e },~~T=\sqrt{\frac{JL}{k_e k_m }},~~\xi =\frac{R\sqrt{J}}{2\sqrt{Lk_e k_m }}$.

Типовое звено - \textit{колебательное 2-го порядка}.

\textbf{Запишем аналитические выражения для временных и частотных характеристик:}

\begin{par}
$$W(j\omega )=\frac{K}{1-T^2 \omega^2 +j(2\xi T\omega )}=\frac{K(1-T^2 \omega^2 -j2\xi T\omega )}{(1-T^2 \omega^2 )^2 +(2\xi T\omega )^2 }=\frac{K(1-T^2 \omega^2 )}{(1-T^2 \omega^2 )^2 +(2\xi T\omega )^2 }-j\frac{2K\xi T\omega }{(1-T^2 \omega^2 )^2 +(2\xi T\omega )^2 }$$
\end{par}

Весовая:

\begin{par}
 $$\begin{array}{l}
w(t)={\mathcal{L}}^{-1} \lbrace W(s)\rbrace ={\mathcal{L}}^{-1} \left\lbrace \frac{K}{T^2 s^2 +2\xi Ts+1}\right\rbrace ={\mathcal{L}}^{-1} \left\lbrace \frac{K/T^2 }{s^2 +2\frac{\xi }{T}s+\frac{1}{T^2 }}\right\rbrace =\\
={\mathcal{L}}^{-1} \left\lbrace \frac{K/T^2 }{{\left(s+\frac{\xi }{T}\right)}^2 +\frac{1-\xi^2 }{T^2 }}\right\rbrace ={\mathcal{L}}^{-1} \left\lbrace \frac{K/T^2 }{{\left(s+\frac{\xi }{T}\right)}^2 +{\left(\frac{\sqrt{1-\xi^2 }}{T}\right)}^2 }\right\rbrace =\\
=\frac{K/T^2 }{\frac{\sqrt{1-\xi^2 }}{T}}e^{-\frac{\xi }{T}t} \sin \left(\frac{\sqrt{1-\xi^2 }}{T}t\right)=\frac{K}{T\sqrt{1-\xi^2 }}e^{-\frac{\xi }{T}t} \sin \left(\frac{\sqrt{1-\xi^2 }}{T}t\right)
\end{array}$$
\end{par}

Переходная:

\begin{par}
 $$\begin{array}{l}
h(t)={\mathcal{L}}^{-1} \left\lbrace \frac{W(s)}{s}\right\rbrace ={\mathcal{L}}^{-1} \left\lbrace \frac{K}{s(T^2 s^2 +2\xi Ts+1)}\right\rbrace ={\mathcal{L}}^{-1} \left\lbrace \frac{K}{s}-\frac{K(s+\frac{2\xi }{T})}{s^2 +\frac{2\xi }{T}s+\frac{1}{T^2 }}\right\rbrace =\\
={\mathcal{L}}^{-1} \left\lbrace \frac{K}{s}-K\cdot \frac{s+\frac{\xi }{T}}{(s+\frac{\xi }{T})^2 +{\left(\frac{\sqrt{1-\xi^2 }}{T}\right)}^2 }-\frac{K\xi }{T}\cdot \frac{1}{(s+\frac{\xi }{T})^2 +{\left(\frac{\sqrt{1-\xi^2 }}{T}\right)}^2 }\right\rbrace =\\
=K-Ke^{-\frac{\xi }{T}t} \left\lbrack \cos \left(\frac{\sqrt{1-\xi^2 }}{T}t\right)+\frac{\xi }{\sqrt{1-\xi^2 }}\sin \left(\frac{\sqrt{1-\xi^2 }}{T}t\right)\right\rbrack 
\end{array}$$
\end{par}

АЧХ: 

\begin{par}
$$\begin{array}{l}
A(\omega )=|W(j\omega )|=\sqrt{{\text{Re}}^2 (W)+{\text{Im}}^2 (W)}=\sqrt{\frac{K^2 }{(1-T^2 \omega^2 )^2 +(2\xi T\omega )^2 }}=\\
=\frac{K}{\sqrt{(1-T^2 \omega^2 )^2 +(2\xi T\omega )^2 }}
\end{array}$$
\end{par}

ФЧХ: $\varphi (\omega )=\arg W(j\omega )=\text{atan2}\left(-2\xi T\omega ,\,1-T^2 \omega^2 \right)=$

ЛАЧХ: $L(\omega )=20\lg (A(\omega ))=20\lg K-10\lg \left\lbrack (1-T^2 \omega^2 )^2 +(2\xi T\omega )^2 \right\rbrack =$

ЛФЧХ: $\varphi (\omega )=-\text{atan}\left(\frac{2\xi T\omega }{1-T^2 \omega^2 }\right)$










\textbackslash{}subsection\{Конденсируй-умножай\}

Дано уравнение конденсатора:

$I=C\frac{dU}{dt}$.

Возьмем из \textbf{Таблицы 2} инструкции по выполнению данной л/р значение, которое соответствует нашему варианту, для следующих величин:

• $C=354$ мкФ — ёмкость конденсатора.

Входом объекта считаем $I(t)$, а выходом $U(t)$.

\textbf{Найдём передаточную функцию и определим тип звена:}

$I=C\dot{U}$, Передаточная функция: $W(s)=\frac{K}{s}$, где $K=\frac{1}{C}$.

Типовое звено - \textit{идеальное интегрирующее.}

\textbf{Запишем аналитические выражения для временных и частотных характеристик:}

\begin{par}
$$W(j\omega )=\frac{K}{j\omega }=-j\frac{K}{\omega }=0-j\frac{K}{\omega }$$
\end{par}

Весовая: $w(t)={\mathcal{L}}^{-1} \lbrace W(s)\rbrace ={\mathcal{L}}^{-1} \left\lbrace \frac{K}{s}\right\rbrace =K$

Переходная: $h(t)={\mathcal{L}}^{-1} \left\lbrace \frac{W(s)}{s}\right\rbrace ={\mathcal{L}}^{-1} \left\lbrace \frac{K}{s^2 }\right\rbrace =Kt$

АЧХ: $A(\omega )=|W(j\omega )|=\sqrt{{\text{Re}}^2 (W)+{\text{Im}}^2 (W)}=\sqrt{0^2 +{\left(\frac{K}{\omega }\right)}^2 }=\frac{K}{\omega }$

ФЧХ: $\varphi (\omega )=\arg W(j\omega )=\text{atan2}\left(-\frac{K}{\omega },\,0\right)=-\frac{\pi }{2}$

ЛАЧХ: $L(\omega )=20\lg (A(\omega ))=20\lg K-20\lg \omega$

ЛФЧХ: $\varphi (\omega )=-\frac{\pi }{2}$








\textbackslash{}subsection\{Пружинка\}

Даны уравнения пружинного маятника:

$F_{\text{упр}} =-k\cdot x$, $F=m\cdot a$

Возьмем из \textbf{Таблицы 3} инструкции по выполнению данной л/р значения, которые соответствуют нашему варианту, для следующих величин:

• $M=26$ кг — масса груза;

• $k=72$ Н/м — коэффициент жесткости пружины.

Входом объекта считаем $F_{ext} (t)$ (некую внешнюю силу, направленную соосно движению маятника), а выходом $x(t)$. Считаем, что маятник движется ортогонально силе тяжести (см. рисунок ???).

???РИСУНОК???

\textbf{Найдём передаточную функцию и определим тип звена:}

\begin{par}
$$F=m\cdot a=m\ddot{x} ,~~m\ddot{x} =-kx+F_{ext} ,~~\frac{m}{k}\ddot{x} +x=\frac{F_{ext} }{k}$$
\end{par}

Передаточная функция: $W(s)=\frac{K}{T^2 s^2 +1}$, где $K=\frac{1}{k},~~T=\sqrt{\frac{m}{k}}$.

Типовое звено - \textit{консервативное}.

\textbf{Запишем аналитические выражения для временных и частотных характеристик:}

\begin{par}
$$W(j\omega )=\frac{K}{1-T^2 \omega^2 }$$
\end{par}

Весовая: $w(t)={\mathcal{L}}^{-1} \lbrace W(s)\rbrace ={\mathcal{L}}^{-1} \left\lbrace \frac{K}{T^2 s^2 +1}\right\rbrace =\frac{K}{T}\sin \left(\frac{t}{T}\right)$

Переходная: $h(t)={\mathcal{L}}^{-1} \left\lbrace \frac{W(s)}{s}\right\rbrace ={\mathcal{L}}^{-1} \left\lbrace \frac{K}{s(T^2 s^2 +1)}\right\rbrace =K\left\lbrack 1-\cos \left(\frac{t}{T}\right)\right\rbrack$

АЧХ: $A(\omega )=|W(j\omega )|=\sqrt{{\text{Re}}^2 (W)+{\text{Im}}^2 (W)}=\frac{K}{\left|1-T^2 \omega^2 \right|}$

ФЧХ: $\varphi (\omega )=\left\lbrace \begin{array}{cc}
0, & \omega <\frac{1}{T},\\
\pi , & \omega >\frac{1}{T},
\end{array}\right.$

ЛАЧХ: $L(\omega )=20\lg A(\omega )=20\lg K-20\lg \left(1-T^2 \omega^2 \right)$

ЛФЧХ: $\varphi (\omega )=\left\lbrace \begin{array}{cc}
0, & \omega <\frac{1}{T},\\
\pi , & \omega >\frac{1}{T},
\end{array}\right.$



